\documentclass[a4paper, serif, 10pt]{moderncv}

%% moderncv
% color and style
\moderncvcolor{black}
\moderncvstyle{banking}

%% page layout/margins
\usepackage[top=2.5cm, bottom=2.5cm, left=1.5cm, right=1.5cm]{geometry}

\nopagenumbers{}

%% Babel config for Spanish language
% \usepackage[spanish]{babel} has some conflicting issues with moderncv
% so we have to use enable the following options to make it work
%
% ref:
% - https://tex.stackexchange.com/a/140161/36007
\usepackage[spanish,es-lcroman]{babel}

%% fontspec
\usepackage{fontspec}

\IfFontExistsTF{Linux Libertine}{
  \setmainfont[Ligatures={TeX, Common}, Numbers=OldStyle]{Linux Libertine}
}{}
\IfFontExistsTF{Linux Libertine O}{
  \setmainfont[Ligatures={TeX, Common}, Numbers=OldStyle]{Linux Libertine O}
}{}

%% CTeX
% CJK support, used to show CJK characters in the resume
%
% - heading=false: disable ctex heading
% - punct=kaiming: use kaiming punctuations styles for CJK
% - scheme=plain: use plain scheme, do not override \normalsize font size
% - space=auto: space settings for CJK characters
%
% CJK fonts are resolved per locale language, e.g. Noto Serif/Sans CJK JP
% for Japanese resumes, with fallback to the next region if not installed
%
% ref:
% - http://ctan.mirrorcatalogs.com/language/chinese/ctex/ctex.pdf
\usepackage[UTF8, heading=false, punct=kaiming, scheme=plain, space=auto]{ctex}

\IfFontExistsTF{Noto Serif CJK SC}{
  \setCJKmainfont{Noto Serif CJK SC}
}{}
\IfFontExistsTF{Noto Sans CJK SC}{
  \setCJKsansfont{Noto Sans CJK SC}
}{}

%% line spacing
\usepackage{setspace}
\setstretch{1.125}

%% URL styling - use normal text instead of monospace
\urlstyle{same}

%% Patch \httplink and \httpslink to handle full URLs
% The original moderncv macros always prepend http:// or https:// to URLs.
% This causes issues when the URL already has a protocol (e.g., https://example.com)
% resulting in malformed URLs like https://https://example.com.
% This patch checks if the URL already contains "://" and uses it directly if so.
% Note: We use \str_set:Nx (x-expansion) to expand macros like \@homepage first.
\ExplSyntaxOn
\str_new:N \g_yamlresume_url_str

\RenewDocumentCommand{\httpslink}{O{}m}{
  \str_set:Nx \g_yamlresume_url_str {#2}
  \str_if_in:NnTF \g_yamlresume_url_str {://}
    { \url{#2} }
    { \url{https://#2} }
}

\RenewDocumentCommand{\httplink}{O{}m}{
  \str_set:Nx \g_yamlresume_url_str {#2}
  \str_if_in:NnTF \g_yamlresume_url_str {://}
    { \url{#2} }
    { \url{http://#2} }
}
\ExplSyntaxOff

\name{Lucía Ferrer Navarro}{}
\title{Analista de Negocio Senior | Transformación digital y mejora de procesos}
\phone[mobile]{+34 612 345 678}
\email{hola@luciaferrer.es}
\homepage{https://www.luciaferrer.es}

\address{Calle de Alcalá 45, 3º B, Madrid, Comunidad de Madrid, España, 28014}{}{}

\extrainfo{{\small \faLinkedin}\ \href{https://www.linkedin.com/in/lucia-ferrer-navarro}{@lucia-ferrer-navarro} {} {} {} • {} {} {} 
{\small \faMedium}\ \href{https://medium.com/@luciaferrer}{@luciaferrer} {} {} {} • {} {} {} 
{\small \faGithub}\ \href{https://github.com/luciaferrer}{@luciaferrer}}

\begin{document}

\maketitle

\section{Información básica}

\cvline{}{Analista de negocio con más de seis años de experiencia en proyectos de transformación digital, optimización de procesos y definición de soluciones basadas en datos.
%
\begin{itemize}
\item Especializada en levantamiento de requisitos, modelado de procesos BPMN y priorización de alcance con stakeholders técnicos y de negocio.
\item Experiencia en entornos ágiles, análisis funcional, definición de KPIs y elaboración de dashboards ejecutivos en Power BI.
\item Capacidad demostrada para traducir necesidades de negocio en historias de usuario, criterios de aceptación y hojas de ruta accionables.
\item Enfoque analítico, comunicación clara y orientación a resultados medibles.
\end{itemize}}
\section{Educación}

\cventry{sept 2018–jul 2020}
        {Maestría, Análisis de Negocio y Transformación Digital, Puntuación: 8,9/10}
        {Universidad Carlos III de Madrid}
        {\url{https://www.uc3m.es}}
        {}
        {\begin{itemize}
\item Formación orientada a la definición de soluciones de negocio, análisis de datos y liderazgo de iniciativas de mejora.
\item Trabajo fin de máster centrado en la priorización de requisitos y la medición de impacto operativo en entornos retail.
\end{itemize}
\textbf{Cursos}: Dirección Estratégica y Modelos de Negocio, Análisis Financiero para la Toma de Decisiones, Gestión de Procesos y BPMN, Analítica de Datos Aplicada, Metodologías Ágiles y Gestión de Producto, Gestión del Cambio y Comunicación, Investigación de Mercados, Gestión de Proyectos}

\cventry{sept 2014–jul 2018}
        {Licenciatura, Administración y Dirección de Empresas, Puntuación: 8,2/10}
        {Universitat de València}
        {\url{https://www.uv.es}}
        {}
        {\begin{itemize}
\item Base sólida en finanzas, operaciones y sistemas de información para el análisis de procesos y la toma de decisiones.
\item Participación en proyectos interdisciplinares de consultoría junior para pequeñas empresas locales.
\end{itemize}
\textbf{Cursos}: Contabilidad Financiera, Estadística Empresarial, Economía de la Empresa, Sistemas de Información, Marketing Estratégico, Derecho Mercantil}
\section{Experiencia laboral}

\cventry{ene 2022 hasta la fecha}
        {Analista de Negocio Senior}
        {Aurea Retail Ibérica}
        {\url{https://www.aurearetail.es}}
        {}
        {\begin{itemize}
\item Lideré el análisis de negocio para la modernización del proceso de alta de clientes, reduciendo el tiempo de tramitación en un 35 \%.
\item Facilité talleres de descubrimiento con negocio, tecnología y operaciones para alinear alcance, prioridades y criterios de aceptación.
\item Definí historias de usuario, reglas de negocio y flujos BPMN para equipos de desarrollo en entornos Scrum y Kanban.
\item Diseñé cuadros de mando en Power BI para el seguimiento de KPIs comerciales y operativos, mejorando la visibilidad de incidencias.
\item Coordiné la gestión del cambio y la documentación funcional para más de 200 usuarios internos.
\end{itemize}
\textbf{Palabras clave}: Análisis de requisitos, Modelado BPMN, Power BI, Scrum, Gestión del cambio}

\cventry{sept 2020–dic 2021}
        {Analista de Negocio}
        {Grupo Meridiano Seguros}
        {\url{https://www.grupomeridianoseguros.es}}
        {}
        {\begin{itemize}
\item Analicé procesos de siniestros y contratación para identificar cuellos de botella y oportunidades de automatización.
\item Elaboré especificaciones funcionales, casos de uso y matrices de trazabilidad para el equipo de desarrollo.
\item Colaboré con negocio y cumplimiento normativo para validar requisitos regulatorios y criterios de aceptación.
\item Implementé informes en SQL y Excel avanzado para el seguimiento de productividad y calidad del servicio.
\end{itemize}
\textbf{Palabras clave}: Análisis funcional, SQL, Casos de uso, Automatización, KPIs}

\cventry{feb 2019–ago 2020}
        {Analista Funcional Junior}
        {Consultora Ítaca}
        {\url{https://www.consultoraitaca.es}}
        {}
        {\begin{itemize}
\item Apoyé la definición de requisitos y la documentación de procesos en proyectos de digitalización para pymes.
\item Participé en reuniones de toma de requisitos, análisis de brechas y preparación de propuestas de solución.
\item Realicé pruebas funcionales y documenté incidencias para asegurar la calidad de las entregas.
\item Preparé manuales de usuario y materiales de formación para la implantación de nuevas herramientas.
\end{itemize}
\textbf{Palabras clave}: Digitalización, Pruebas funcionales, Documentación funcional, Análisis de brechas, Formación}
\section{Idiomas}

\cvline{Español}{Competencia nativa o bilingüe \hfill \textbf{Palabras clave}: Lengua materna}
\cvline{Inglés}{Competencia profesional plena \hfill \textbf{Palabras clave}: C1 Cambridge Advanced}
\cvline{Francés}{Competencia limitada de trabajo \hfill \textbf{Palabras clave}: A2 DELF}
\section{Competencias}

\cvline{Análisis de Negocio}{Experto \hfill \textbf{Palabras clave}: Análisis de requisitos, Modelado BPMN, Historias de usuario, Casos de uso, Gestión de stakeholders}
\cvline{Datos y Reporting}{Avanzado \hfill \textbf{Palabras clave}: SQL, Power BI, Tableau, Excel avanzado, KPIs}
\cvline{Herramientas}{Avanzado \hfill \textbf{Palabras clave}: Jira, Confluence, Azure DevOps, Miro, Visio}
\cvline{Metodologías}{Avanzado \hfill \textbf{Palabras clave}: Scrum, Kanban, Lean Six Sigma, Design Thinking, Gestión del cambio}
\section{Premios}

\cventry{nov 2023}
        {Aurea Retail Ibérica}
        {Premio a la Excelencia Operacional}
        {}
        {}
        {Reconocimiento por la mejora del proceso de alta de clientes y la reducción de tiempos de tramitación mediante análisis de negocio y automatización.}

\cventry{may 2022}
        {Retail Data Summit}
        {Finalista del Datathon de Retail}
        {}
        {}
        {Proyecto finalista centrado en la detección de patrones de abandono y la priorización de acciones comerciales a partir de datos transaccionales.}
\section{Certificados}

\cventry{mar 2024}
        {IIBA}
        {Certificación CBAP en Análisis de Negocio}
        {\url{https://www.iiba.org/certification/cbap/}}
        {}
        {}

\cventry{jun 2022}
        {Scrum.org}
        {Scrum Master Profesional I}
        {\url{https://www.scrum.org/professional-scrum-master-i-certification}}
        {}
        {}

\cventry{sept 2023}
        {Microsoft}
        {Analista de Datos Power BI Associate}
        {\url{https://learn.microsoft.com/en-us/certifications/power-bi-data-analyst-associate/}}
        {}
        {}
\section{Publicaciones}

\cventry{mar 2024}
        {Guía práctica para priorizar requisitos con MoSCoW}
        {Blog de Producto y Datos}
        {\url{https://www.productoydatos.es/guias/moscow-requisitos}}
        {}
        {\begin{itemize}
\item Artículo sobre técnicas de priorización de requisitos en proyectos con múltiples stakeholders.
\item Incluye criterios prácticos para equilibrar valor de negocio, esfuerzo técnico y riesgo.
\end{itemize}}
\section{Proyectos}

\cventry{mar 2023–nov 2023}
        {Proyecto de transformación operativa para simplificar la contratación y mejorar la experiencia del cliente.}
        {Rediseño del proceso de alta de clientes}
        {\url{https://www.aurearetail.es/casos/alta-clientes}}
        {}
        {\begin{itemize}
\item Mapeo del proceso actual y definición del proceso objetivo con negocio, operaciones y tecnología.
\item Especificación de requisitos funcionales, reglas de validación y criterios de aceptación.
\item Seguimiento de indicadores de adopción, tiempo de tramitación y satisfacción de usuarios internos.
\end{itemize}
\textbf{Palabras clave}: BPMN, Power BI, Gestión del cambio, Experiencia de cliente}

\cventry{ene 2022–dic 2022}
        {Solución de reporting para unificar métricas de ventas, tráfico y conversión en tienda y canal online.}
        {Cuadro de mando comercial omnicanal}
        {\url{https://www.aurearetail.es/casos/cuadro-mando-omnicanal}}
        {}
        {\begin{itemize}
\item Definición de KPIs y reglas de negocio junto a dirección comercial y equipos de datos.
\item Diseño de dashboards en Power BI para el seguimiento semanal de objetivos por canal y región.
\item Documentación de fuentes de datos, definiciones métricas y procesos de actualización.
\end{itemize}
\textbf{Palabras clave}: Power BI, KPIs, Omnicanalidad, Modelo de datos}
\section{Intereses}

\cvline{Analítica y visualización de datos}{Power BI, Storytelling con datos, Datos abiertos}
\cvline{Sostenibilidad}{Economía circular, Movilidad urbana}
\cvline{Lectura}{Ensayo, Novela negra}
\section{Voluntariado}

\cventry{sept 2022 hasta la fecha}
        {Mentora de jóvenes analistas}
        {Asociación de Analistas de Negocio de España}
        {\url{https://www.aanes.es}}
        {}
        {\begin{itemize}
\item Acompaño a profesionales junior en la preparación de entrevistas, elaboración de currículos y desarrollo de competencias de análisis.
\item Participo en sesiones formativas sobre requisitos, modelado de procesos y comunicación con stakeholders.
\end{itemize}}
    
\end{document}